% ===========================================================================
% 8 Limitations and future work, 9 Conclusion
% Drafted with AI assistance (Claude). Rewrite in your own words before submitting.
% Revised on 4 Nov 2026 after the supervisor's review: limitations kept short.
% ===========================================================================
\section{Limitations and Future Work}
\label{sec:limitations}

The method was developed on SciFact and validated on one further biomedical dataset, so larger
corpora and other domains remain to be tested, and the failure categories, assigned by a single
annotator with AI assistance, would benefit from independent validation. Future work includes
selecting the fusion weight for each claim from inexpensive signals and fine-tuning on larger
biomedical corpora, to capture relations such as drug-to-class terminology that a few hundred
claims cannot teach.

\section{Conclusion}
\label{sec:conclusion}

This study addressed one part of scientific question answering, the retrieval of evidence for a
claim, and improved it. Fine-tuning a 22-million-parameter encoder on the domain's training
claims and fusing it with keyword search at a weight below the break-even point of
Proposition~1 retrieved evidence more accurately than the standard hybrid pipeline with a
cross-encoder reranker (nDCG@10 of 0.759 against 0.715 on SciFact, $p = 0.003$), at about a
third of its time per claim. Proposition~1 explains why the weight works: below $61/160$, keyword
search can only re-order the candidates of the dense retriever. On NFCorpus, the same method
matched the reranked pipeline at a sixth of its cost. Because literature assistants depend on
this retrieval step, the results show how it can be made both more accurate and less expensive
with an openly evaluated method that runs on a laptop.
